\documentclass[11pt,a4paper]{article}

\usepackage[utf8]{inputenc}
\usepackage[T1]{fontenc}
\usepackage[spanish,es-nodecimaldot]{babel}
\usepackage{lmodern}
\usepackage{geometry}
\usepackage{xcolor}
\usepackage{hyperref}
\usepackage{enumitem}
\usepackage{listings}
\usepackage{longtable}
\usepackage{tabularx}
\usepackage{booktabs}
\usepackage{fancyhdr}
\usepackage{microtype}

\geometry{top=2.1cm,bottom=2.1cm,left=2.2cm,right=2.2cm}
\hypersetup{
  colorlinks=true,
  linkcolor=blue!55!black,
  urlcolor=blue!60!black,
  pdftitle={Guía detallada de la implementación CI/CD para Azure Databricks},
  pdfauthor={Proyecto Databricks CI/CD}
}

\definecolor{dbblue}{HTML}{1B4F72}
\definecolor{dbcyan}{HTML}{117A8B}
\definecolor{lightblue}{HTML}{EAF2F8}
\definecolor{codebg}{HTML}{F6F8FA}
\definecolor{codegray}{HTML}{586069}

\pagestyle{fancy}
\fancyhf{}
\lhead{Proyecto Azure Databricks CI/CD}
\rhead{Guía de aprendizaje}
\cfoot{\thepage}
\setlength{\headheight}{14pt}

\setlist[itemize]{leftmargin=1.4em,itemsep=0.25em,topsep=0.4em}
\setlist[enumerate]{leftmargin=1.6em,itemsep=0.35em,topsep=0.4em}
\setlength{\parindent}{0pt}
\setlength{\parskip}{0.55em}

\lstdefinestyle{projectcode}{
  backgroundcolor=\color{codebg},
  basicstyle=\ttfamily\small,
  breaklines=true,
  columns=fullflexible,
  commentstyle=\color{codegray},
  frame=single,
  framerule=0.3pt,
  rulecolor=\color{black!20},
  showstringspaces=false,
  tabsize=2
}
\lstset{style=projectcode}

\newcommand{\concepto}[1]{\textcolor{dbblue}{\textbf{#1}}}
\newcommand{\estado}[1]{\colorbox{lightblue}{\strut\textbf{#1}}}

\begin{document}

\begin{titlepage}
  \centering
  \vspace*{2.2cm}
  {\color{dbblue}\rule{\textwidth}{1.2pt}}\\[1cm]
  {\Huge\bfseries Guía detallada de la implementación\\[0.3cm]
  CI/CD para Azure Databricks\par}
  \vspace{0.7cm}
  {\Large GitHub Actions, Terraform, seguridad y Declarative Automation Bundles\par}
  \vspace{1cm}
  {\color{dbblue}\rule{\textwidth}{1.2pt}}\\[1.5cm]
  {\large Documento didáctico del estado actual del repositorio\par}
  \vfill
  \begin{tabular}{rl}
    Proyecto: & \texttt{proyecto-databricks-cicd} \\
    Alcance actual: & Entorno Dev preparado, sin recursos cloud creados \\
    Fecha: & 8 de octubre de 2026 \\
    Presupuesto objetivo: & S/200 mensuales \\
  \end{tabular}
  \vfill
  {\small Este documento explica qué se creó, por qué existe cada pieza, cómo se validó y qué falta ejecutar.}
\end{titlepage}

\tableofcontents
\newpage

\section{Propósito y estado actual}

El proyecto es un laboratorio profesional y progresivo para aprender a llevar código Python, SQL y
notebooks desde GitHub hasta Azure Databricks. La primera meta no es desplegar de inmediato en tres
entornos, sino demostrar de punta a punta el flujo hacia \textbf{Dev}, controlar los costos y recién
después incorporar QA y Prod.

\subsection{Qué está implementado en el repositorio}

\begin{itemize}
  \item Estructura de código Python, SQL, notebooks y pruebas.
  \item Configuración reproducible de herramientas de calidad y seguridad.
  \item Terraform para crear el workspace Azure Databricks Dev y su presupuesto.
  \item Terraform para configurar identidad OIDC, catálogo, esquema y permisos del workspace.
  \item Un Bundle que publica una carpeta normal y administra el Job de smoke test.
  \item Un workflow de validación de Pull Requests.
  \item Un workflow de despliegue automático hacia Dev después del merge a \texttt{main}.
  \item Documentación operativa, reglas de contribución y registros de decisiones arquitectónicas.
\end{itemize}

\subsection{Qué todavía no se ha ejecutado}

\begin{itemize}
  \item No existe todavía una suscripción Azure asociada al proyecto.
  \item No se ejecutó \texttt{terraform plan} con valores reales ni \texttt{terraform apply}.
  \item No se creó ningún workspace Databricks y no existe consumo cloud generado por el proyecto.
  \item No se creó todavía el repositorio remoto de GitHub ni el GitHub Environment \texttt{dev}.
  \item QA y Prod están diseñados conceptualmente, pero no se aprovisionan en esta fase.
\end{itemize}

Esta distinción es importante: \concepto{código de infraestructura} no equivale a
\concepto{infraestructura creada}. Terraform sólo realiza cambios reales cuando se revisa y ejecuta
un \texttt{apply} autorizado.

\subsection{Qué se hace manualmente y qué automatiza Terraform}

La persona propietaria crea la cuenta Microsoft/Azure, activa MFA, registra facturación, convierte la
suscripción gratuita a Pay-As-You-Go, acepta términos, selecciona la suscripción mediante Azure CLI y
autoriza cada \texttt{apply}. Son decisiones de identidad y facturación que no deben delegarse a un
pipeline.

Después de ese bootstrap, Terraform crea el Resource Group, presupuesto, workspace Dev, identidad y
permisos. No se deben crear manualmente esos recursos en el portal. Terraform tampoco despliega cada
notebook: los Bundles administran el código y Jobs con una cadencia mucho mayor que la infraestructura.

\section{Arquitectura y flujo general}

El flujo elegido combina integración continua sobre \texttt{main} con promoción de commits
inmutables. \texttt{main} representa código integrado; no representa automáticamente la versión de
Producción.

\begin{lstlisting}
feature/* o fix/*
        |
        | Pull Request: calidad, pruebas y seguridad
        v
      main  ---- despliegue automatico ----> Dev
        |
        | tag vX.Y.Z-rc.N
        v
       QA   ---- pruebas + aprobacion ----> mismo SHA
        |
        | tag final vX.Y.Z
        v
      Prod
\end{lstlisting}

Si un release candidate falla en QA, Prod continúa usando la versión estable anterior. La corrección
se hace en una nueva rama \texttt{fix/*}, se integra mediante otro PR y genera un nuevo SHA y un nuevo
candidato. No se mueve ni se reutiliza el tag fallido.

\subsection{Por qué se promueve el mismo SHA}

Un \concepto{SHA} identifica de forma única un commit Git. Promover el mismo SHA reduce la pregunta
``¿QA probó exactamente lo que llegó a Prod?'' a una comparación verificable. No se recompila ni se
edita el código durante la promoción.

\section{Mapa del repositorio}

\begin{longtable}{p{0.34\textwidth}p{0.58\textwidth}}
\toprule
\textbf{Ruta} & \textbf{Responsabilidad} \\
\midrule
\endhead
\texttt{src/databricks\_cicd/} & Código Python reutilizable y contrato de configuración por entorno. \\
\texttt{tests/} & Pruebas unitarias ejecutadas con Pytest. \\
\texttt{sql/} & SQL versionado y validado con dialecto Databricks. \\
\texttt{proyecto\_test/} & Árbol de notebooks y configuración reflejado en el workspace. \\
\texttt{resources/} & Jobs y otros recursos declarados por el Bundle. \\
\texttt{databricks.yml} & Targets, rutas y variables del Bundle. \\
\texttt{scripts/} & Utilidades del pipeline, como validación de variables requeridas. \\
\texttt{infra/azure/dev/} & Recursos Azure: presupuesto, resource group y workspace Dev. \\
\texttt{infra/databricks/dev/} & Configuración interna de Databricks después de crear el workspace. \\
\texttt{infra/modules/} & Módulos Terraform reutilizables. \\
\texttt{.github/workflows/} & CI de Pull Requests y despliegue Dev. \\
\texttt{docs/adr/} & Decisiones de arquitectura y sus motivos. \\
\texttt{pyproject.toml} & Metadatos Python, dependencias de desarrollo y configuración de herramientas. \\
\texttt{.gitignore} & Evita versionar state, variables locales, entornos virtuales y credenciales. \\
\bottomrule
\end{longtable}

\section{Python, entorno virtual y dependencias}

\subsection{Qué es un entorno virtual}

Un \concepto{entorno virtual} o \texttt{venv} es una carpeta aislada que contiene un intérprete de
Python y sus paquetes. Evita mezclar dependencias del proyecto con paquetes globales del equipo.

Durante la implementación se creó localmente \texttt{.venv/} y se instalaron las dependencias de
desarrollo declaradas en \texttt{pyproject.toml}. La carpeta está ignorada por Git porque:

\begin{itemize}
  \item puede reconstruirse a partir de la configuración;
  \item cambia entre Windows, Linux y macOS;
  \item versionarla añadiría miles de archivos ajenos al código del proyecto.
\end{itemize}

El proceso conceptual fue:

\begin{lstlisting}[language=bash]
python -m venv .venv
.venv/Scripts/python -m pip install --editable ".[dev]"
\end{lstlisting}

\texttt{--editable} instala el proyecto apuntando a la carpeta de trabajo. Así, los cambios en
\texttt{src/} se reflejan sin reinstalar el paquete. El extra \texttt{[dev]} instala únicamente las
herramientas necesarias para desarrollar y validar.

\subsection{Qué contiene pyproject.toml}

\texttt{pyproject.toml} centraliza:

\begin{itemize}
  \item nombre, versión y versión mínima de Python;
  \item sistema de construcción basado en Setuptools;
  \item dependencias de desarrollo;
  \item configuración de Pytest, Ruff y Bandit;
  \item ubicación del paquete dentro de \texttt{src/}.
\end{itemize}

Las dependencias de validación incluyen Ruff, Pytest, SQLFluff, Yamllint, Bandit,
\texttt{detect-secrets} y \texttt{pip-audit}. No son servicios cloud: son programas ejecutados en el
equipo local o en un runner de GitHub Actions.

\section{Validaciones de código y seguridad}

\begin{longtable}{p{0.20\textwidth}p{0.35\textwidth}p{0.36\textwidth}}
\toprule
\textbf{Herramienta} & \textbf{Qué valida} & \textbf{Qué evita} \\
\midrule
\endhead
Ruff & Estilo, imports, errores Python y formato. & Código inconsistente, imports sin usar y defectos simples. \\
Pytest & Comportamiento esperado mediante pruebas. & Cambios que rompen reglas ya verificadas. \\
SQLFluff & Sintaxis y estilo SQL con dialecto Databricks. & SQL difícil de revisar o incompatible con las reglas elegidas. \\
Yamllint & Estructura y estilo de archivos YAML. & Workflows mal indentados o configuración ambigua. \\
Bandit & Análisis estático de seguridad para Python. & Patrones inseguros conocidos dentro del código. \\
detect-secrets & Patrones y entropía compatibles con secretos. & Tokens o claves enviados accidentalmente a Git. \\
pip-audit & Vulnerabilidades publicadas de dependencias. & Usar versiones con CVE o avisos de seguridad conocidos. \\
Terraform validate & Sintaxis, tipos y relaciones de configuración. & Planes inválidos antes de contactar Azure. \\
\bottomrule
\end{longtable}

\subsection{Pruebas implementadas}

El módulo \texttt{config.py} modela tres valores obligatorios: entorno, catálogo y esquema. Las pruebas
verifican que:

\begin{enumerate}
  \item una configuración Dev válida se normalice correctamente;
  \item una variable ausente cause un error explícito;
  \item un entorno distinto de Dev, QA o Prod sea rechazado.
\end{enumerate}

Este patrón mueve errores de configuración hacia el inicio del pipeline. Es preferible fallar antes
de autenticarse o modificar Databricks que descubrir el problema a mitad del despliegue.

\subsection{Qué ocurrió durante la validación real}

La primera ejecución no fue artificialmente perfecta. Encontró:

\begin{itemize}
  \item imports Python sin normalizar;
  \item archivos que Ruff deseaba reformatear;
  \item Yamllint intentando revisar el contenido de \texttt{.venv};
  \item un salto de línea SQL incorrecto;
  \item versiones de \texttt{pip}, Pytest y SQLFluff reportadas con vulnerabilidades conocidas.
\end{itemize}

Las correcciones fueron: ordenar y formatear Python, excluir \texttt{.venv} del lint YAML, corregir el
SQL, actualizar los rangos de dependencias y repetir todas las validaciones. El resultado final fue:

\begin{itemize}
  \item tres pruebas unitarias aprobadas;
  \item Ruff, SQLFluff y Yamllint aprobados;
  \item Bandit sin incidencias;
  \item detección de secretos aprobada;
  \item \texttt{pip-audit} sin vulnerabilidades conocidas en las dependencias auditables.
\end{itemize}

Esto ilustra el propósito de CI: convertir problemas reproducibles en feedback antes del merge.

\section{Terraform desde cero}

\subsection{Conceptos esenciales}

\begin{description}[style=nextline,leftmargin=0pt]
  \item[Provider] Plugin que traduce Terraform a la API de una plataforma. Aquí se usan AzureRM y Databricks.
  \item[Resource] Objeto que Terraform crea o administra, por ejemplo un resource group o catálogo.
  \item[Variable] Entrada configurable; evita codificar IDs o valores de entorno dentro de los recursos.
  \item[Output] Dato producido por una capa para uso humano o para configurar la siguiente.
  \item[Module] Conjunto reutilizable de recursos con entradas y salidas.
  \item[State] Registro que vincula la configuración con recursos reales. Puede contener datos sensibles.
  \item[Lockfile] Archivo que fija las versiones exactas de providers y sus hashes.
\end{description}

\subsection{Ciclo de comandos}

\begin{lstlisting}[language=bash]
terraform init
terraform fmt -check
terraform validate
terraform plan -out=tfplan
terraform apply tfplan
\end{lstlisting}

\texttt{init} prepara providers; \texttt{fmt} revisa formato; \texttt{validate} comprueba la
configuración; \texttt{plan} muestra cambios propuestos; \texttt{apply} ejecuta exactamente el plan.
Sólo los dos últimos requieren valores y permisos cloud reales, y sólo \texttt{apply} crea o modifica
recursos.

\subsection{Por qué existen dos capas}

\textbf{Capa Azure:} crea el contenedor cloud y el workspace. Incluye el presupuesto y etiquetas de
costos. Usa el provider AzureRM.

\textbf{Capa Databricks:} se ejecuta cuando el workspace ya responde. Crea la identidad Databricks,
política OIDC, permisos, catálogo, esquema, carpeta del proyecto y ruta protegida para el estado del
Bundle. Usa providers Databricks a nivel de cuenta y workspace.

Separarlas evita intentar configurar una API que todavía no existe y permite destruir primero los
objetos internos, y luego el workspace.

\subsection{State local y archivos ignorados}

En esta fase el state comienza local para facilitar el aprendizaje. Los patrones \texttt{*.tfstate},
\texttt{*.tfvars} y \texttt{*.tfplan} están ignorados. Los ejemplos
\texttt{terraform.tfvars.example} sí se versionan porque sólo contienen marcadores ficticios.

Antes de añadir QA, el state se migrará a Azure Storage con cifrado, control de acceso, versionado y
bloqueo. El state nunca debe pegarse en tickets ni enviarse por correo.

\subsection{Validación efectuada}

Se descargó temporalmente Terraform 1.13.4 con verificación SHA-256 para validar compatibilidad. Se
ejecutaron \texttt{fmt}, \texttt{init -backend=false} y \texttt{validate} sobre ambas capas. Los
providers quedaron fijados por lockfiles. El workflow de CI solicita Terraform 1.16.5.

\texttt{-backend=false} significa que la validación no intentó configurar state remoto. No se ejecutó
\texttt{plan} ni \texttt{apply}; por tanto, no hubo cambios en Azure.

\section{Infraestructura Azure definida}

La capa Azure declara:

\begin{itemize}
  \item un resource group para Dev;
  \item un workspace Azure Databricks Premium;
  \item un managed resource group con nombre controlado;
  \item etiquetas \texttt{project}, \texttt{environment}, \texttt{owner} y \texttt{managed\_by};
  \item un presupuesto mensual con alertas al 50, 75, 90 y 100 por ciento.
\end{itemize}

El módulo restringe intencionalmente \texttt{environment} a \texttt{dev}. Esto hace que intentar crear
QA o Prod en la fase equivocada falle durante la validación.

Azure Budgets alerta, pero no corta el consumo. Además, el importe está expresado en la moneda de
facturación de la suscripción, no necesariamente en soles. Antes del primer \texttt{apply} se debe
convertir S/200 con margen conservador y confirmar la región.

\section{Configuración Databricks definida}

La segunda capa declara:

\begin{itemize}
  \item un service principal exclusivo para CI/CD Dev;
  \item asignación del principal al workspace con rol de usuario, no administrador;
  \item una política de federación OIDC limitada al repositorio y GitHub Environment \texttt{dev};
  \item la carpeta normal \texttt{/Workspace/data/proyecto\_test};
  \item la ruta protegida \texttt{/Workspace/.cicd/proyecto\_test/dev};
  \item \texttt{CAN\_MANAGE} para CI/CD y edición temporal para desarrolladores en Dev;
  \item catálogo \texttt{proyecto\_dev} y esquema \texttt{cicd};
  \item permisos de datos específicos para la identidad del pipeline.
\end{itemize}

El catálogo y el esquema usan \texttt{force\_destroy = false}. Si contienen datos, Terraform debe
negarse a eliminarlos automáticamente. Es una barrera contra destrucción accidental.

\section{GitHub Actions: validación del Pull Request}

El workflow \texttt{validate-pr.yml} se dispara ante un PR hacia \texttt{main} o manualmente. Declara
\texttt{contents: read} y no recibe credenciales cloud.

\subsection{Jobs de CI}

\begin{enumerate}
  \item \textbf{Python, SQL and YAML}: instala el proyecto y ejecuta Ruff, SQLFluff, Yamllint y Pytest.
  \item \textbf{Security checks}: ejecuta detección de secretos, Bandit y \texttt{pip-audit}.
  \item \textbf{Terraform static validation}: revisa formato, inicializa sin backend y valida ambas capas.
\end{enumerate}

Los jobs usan \texttt{timeout-minutes}; la concurrencia cancela validaciones antiguas del mismo PR.
Las acciones principales están fijadas a SHA completos con un comentario que indica la versión legible.

\subsection{Por qué el PR no posee secretos}

Un contribuidor podría modificar el código que ejecuta CI. Entregar credenciales a ese proceso
permitiría exfiltrarlas. Por eso la validación de PR sólo analiza código y no usa \texttt{id-token:
write}, tokens Databricks ni secretos de Azure.

\section{GitHub Actions: despliegue Dev}

\texttt{deploy-dev.yml} se ejecuta después de un push a \texttt{main}, normalmente producido por un
merge aprobado. Usa el GitHub Environment \texttt{dev} y declara sólo:

\begin{lstlisting}
permissions:
  contents: read
  id-token: write
\end{lstlisting}

\texttt{id-token: write} no da permisos generales de escritura en GitHub. Permite pedir un token OIDC
temporal que Databricks validará contra la política configurada.

\subsection{Secuencia exacta del job}

\begin{enumerate}
  \item Descarga exactamente \texttt{github.sha}.
  \item Valida variables obligatorias sin imprimir sus valores.
  \item Instala la CLI oficial con \texttt{databricks/setup-cli}.
  \item Comprueba la identidad OIDC mediante \texttt{databricks current-user me}.
  \item Valida y despliega el Bundle parametrizado con \texttt{github.sha}.
  \item Reconcilia rutas, tipos y hashes contra un manifiesto del repositorio.
  \item Elimina archivos manuales y restaura contenido modificado o eliminado.
  \item Ejecuta el Job serverless de humo declarado por el Bundle.
  \item Escribe un resumen con commit, ruta y entorno, sin secretos.
\end{enumerate}

La concurrencia \texttt{databricks-dev} impide dos despliegues simultáneos. El Bundle conserva una
identidad de despliegue y una segunda ejecución del mismo SHA debe producir el mismo estado.

\section{OIDC y mínimo privilegio}

\concepto{OIDC} permite que GitHub demuestre ante Databricks qué repositorio y environment ejecutan el
job. Databricks intercambia esa prueba por una credencial temporal. No existe un PAT de larga duración
guardado en GitHub.

La política acepta un subject similar a:

\begin{lstlisting}
repo:OWNER/REPOSITORY:environment:dev
\end{lstlisting}

Si el repositorio se renombra, transfiere o intenta usar otro environment, la autenticación falla. Es
un control positivo, no un inconveniente.

El service principal no es workspace admin. Sólo recibe acceso al workspace, a la carpeta del
proyecto, al estado del Bundle y al catálogo necesarios. QA y Prod tendrán identidades diferentes.

\section{Bundle, carpeta normal y smoke test}

El árbol local \texttt{proyecto\_test/} se publica bajo
\texttt{/Workspace/data/proyecto\_test}. No aparece como Git Folder. El archivo
\texttt{databricks.yml} define targets y rutas; \texttt{resources/} define Jobs. El estado técnico
vive fuera de la carpeta editable, bajo \texttt{/Workspace/.cicd}.

Un manifiesto calcula la ruta remota, tipo y SHA-256 de cada archivo. Después de desplegar, el pipeline
elimina objetos desconocidos y compara el contenido exportado. Sus guardas rechazan rutas amplias como
\texttt{/Workspace} o \texttt{/Workspace/data}; sólo puede limpiar la carpeta exacta del proyecto.

El smoke test es deliberadamente pequeño:

\begin{lstlisting}[language=Python]
result = spark.sql("SELECT 1 AS pipeline_is_healthy").first()
assert result.pipeline_is_healthy == 1
\end{lstlisting}

No prueba toda la lógica del proyecto. Confirma que autenticación, Bundle, reconciliación, lectura del
notebook y cómputo serverless funcionan juntos. Las pruebas de negocio crecerán después sin usar datos
reales.

\subsection{Debugging manual en Dev y QA}

Los desarrolladores pueden modificar archivos y notebooks en Dev y QA para diagnosticar. Esos cambios
son temporales: deben copiarse a una rama y pasar por PR. El siguiente despliegue restaura lo que
existe en Git, recrea eliminados y borra agregados manualmente. No existe sincronización automática
workspace a Git. Prod permanece sin edición manual para desarrolladores.

\section{Flujo completo de un cambio}

\begin{enumerate}
  \item El desarrollador crea \texttt{feature/descripcion} desde \texttt{main}.
  \item Trabaja localmente dentro de \texttt{.venv} y ejecuta linters y pruebas.
  \item Abre un PR a \texttt{main}.
  \item GitHub ejecuta calidad, seguridad y Terraform sin credenciales cloud.
  \item Una revisión humana aprueba y se realiza squash merge.
  \item El push resultante a \texttt{main} inicia el workflow Dev.
  \item OIDC entrega una identidad temporal de mínimo privilegio.
  \item Databricks despliega el Bundle y reconcilia la carpeta contra el SHA.
  \item Se ejecuta el smoke test serverless.
  \item El resultado y el SHA quedan registrados en GitHub Actions.
\end{enumerate}

En la fase futura, un tag \texttt{v0.1.0-rc.1} seleccionará un SHA para QA. Si pasa, el tag final
\texttt{v0.1.0} apuntará al mismo SHA y será elegible para Prod.

\section{Controles de costo}

La estrategia para acercarse a S/200 mensuales es:

\begin{itemize}
  \item crear sólo Dev al inicio;
  \item usar serverless o cómputo efímero;
  \item no mantener clústeres clásicos permanentes;
  \item revisar costos y forecast después de cada práctica;
  \item usar alertas progresivas;
  \item destruir el laboratorio si aparece gasto inesperado;
  \item crear QA y Prod sólo cuando exista una prueba que justifique su costo.
\end{itemize}

El orden seguro de destrucción es configuración Databricks primero y capa Azure después. Siempre se
revisa \texttt{terraform plan -destroy} antes de aplicar la destrucción.

\section{Archivos de seguridad y reproducibilidad}

\subsection{.gitignore}

Excluye \texttt{.venv}, caches, archivos \texttt{.env}, certificados, \texttt{terraform.tfvars},
planes y states. Mantiene versionados los ejemplos sin credenciales.

\subsection{Lockfiles Terraform}

Los archivos \texttt{.terraform.lock.hcl} registran provider, versión seleccionada y hashes. Sí deben
versionarse: permiten que CI y otros equipos utilicen la misma versión validada.

\subsection{Dependabot}

Revisa semanalmente dependencias Python, GitHub Actions y providers Terraform. Una actualización no se
aplica directamente: genera un PR que debe pasar por los mismos checks.

\section{Qué significa ``listo'' en este punto}

\estado{Listo localmente}

\begin{itemize}
  \item código y configuración escritos;
  \item pruebas y linters aprobados;
  \item dependencias auditadas;
  \item Terraform inicializado y validado sin backend;
  \item documentación y ADRs disponibles.
\end{itemize}

\estado{Pendiente de credenciales y autorización humana}

\begin{itemize}
  \item alta Pay-As-You-Go y MFA;
  \item elección de región y conversión del presupuesto a la moneda de facturación;
  \item creación del remoto público GitHub;
  \item primer plan Azure con valores reales;
  \item revisión conjunta y primer apply;
  \item configuración de variables del GitHub Environment Dev;
  \item primera ejecución real del pipeline.
\end{itemize}

\section{Glosario breve}

\begin{description}[style=nextline,leftmargin=0pt]
  \item[CI] Integración continua: validación automática de cada cambio.
  \item[CD] Entrega o despliegue continuo: promoción automatizada y controlada.
  \item[Runner] Máquina efímera donde GitHub ejecuta un job.
  \item[Workflow] Archivo YAML que define eventos y jobs de GitHub Actions.
  \item[Environment] Frontera GitHub para variables, secretos y aprobaciones de un entorno.
  \item[Service principal] Identidad no humana usada por automatización.
  \item[OIDC] Federación que proporciona credenciales temporales sin secreto persistente.
  \item[SAST] Análisis estático de seguridad del código.
  \item[SCA] Análisis de vulnerabilidades de dependencias.
  \item[Idempotencia] Repetir una operación produce el mismo estado, sin duplicados.
  \item[Smoke test] Prueba pequeña que confirma que el camino crítico funciona.
  \item[Unity Catalog] Capa de gobierno y permisos para datos y objetos Databricks.
  \item[Release candidate] Versión candidata para QA, todavía no productiva.
\end{description}

\section{Próximo recorrido recomendado}

\begin{enumerate}
  \item Leer \texttt{docs/01-azure-account-and-cost.md} y crear la suscripción personalmente.
  \item Confirmar moneda de facturación, presupuesto equivalente y región disponible.
  \item Crear el repositorio GitHub público y publicar el commit inicial.
  \item Instalar Azure CLI y Terraform desde fuentes oficiales.
  \item Completar \texttt{terraform.tfvars} local sin versionarlo.
  \item Ejecutar juntos \texttt{terraform plan}; explicar cada alta, cambio o destrucción.
  \item Sólo después de revisar el plan, autorizar \texttt{terraform apply} para Dev.
  \item Configurar Databricks y GitHub Environment Dev.
  \item Abrir un PR de prueba y observar cada job.
  \item Medir costo real antes de decidir la creación de QA.
\end{enumerate}

\vfill
\begin{center}
  \color{dbcyan}\textbf{Principio rector: entender cada cambio antes de automatizarlo.}
\end{center}

\end{document}
